\documentclass{situationist-preprint}

\title{Divine Agency and the Limits of the Main Result}
\msid{REL-01-e}
\msfamily{REL-01}
\msdiscipline{Philosophy of Religion}
\mstype{Consequence}
\mskeywords{divine agency; ontological argument; omnipotence; creation; providence; limits of proof}

\begin{document}
\maketitle

\begin{abstract}
The main result of \emph{On the Existence of God} (REL-01-a) is that a maximally great being exists necessarily, and hence actually. We derive what follows from it about divine agency, and find that it is nothing. Maximal excellence is a conjunction of powers, knowledge and moral perfection; each is compatible with every pattern of exercise, including none. We show that any model of the main argument can be extended both by a causal relation in which the maximally great being causes every contingent being and by one in which it causes nothing, without disturbing any premise. Creation, providence, miracle and answered prayer are therefore not consequences of the main result, and arguments for them need premises of another kind.
\end{abstract}

\section{The main result and its vocabulary}

The main argument of REL-01-a concludes that some being is maximally great at the actual world (Theorem~3.4) and therefore exists there and is maximally excellent (Corollary~3.5). Its vocabulary is fixed by Definitions~2.2 to~2.4: a being is maximally excellent at a world if it is omnipotent, omniscient and morally perfect there, and maximally great if it exists and is maximally excellent at every accessible world. Its premises are the S5 conditions (Assumption~2.1) and the possibility premise (Assumption~3.1). No term of the argument denotes an act, and no premise mentions one. REL-01-a accordingly says that the question of divine agency is left to Volume~IX, where the main result is shown to leave the causal role of the maximally great being undetermined. The present note gives the elementary form of that point.

\section{The independence of agency}

Let a \emph{model} of the main argument be a structure $\langle W, @, R, B, E, \mathrm{ME}\rangle$ satisfying Assumptions~2.1 and~3.1 of REL-01-a. Let a \emph{causal extension} of a model add a relation $C(x, y, w)$, read: $x$ causes $y$ to exist at $w$, subject only to factivity ($C(x,y,w)$ implies $E(x,w)$ and $E(y,w)$).

\begin{assumption}[Powers are not exercises]\label{a:powers}
Whether a being is omnipotent, omniscient or morally perfect at a world does not depend on which of its powers it exercises there.
\end{assumption}

Assumption~\ref{a:powers} is the standard reading of the attributes. Omnipotence is defined by what an agent \emph{can} do, both in the stratified sense of REL-04-a (Definition~2.3) and in the sense used in REL-02-a (Definition~2.7); omniscience is a matter of what is known; and moral perfection, as REL-02-a defines it (Definition~2.9), excludes wrong options without selecting among permissible ones. \citet{Pike1969} made the parallel point for divine impeccability: that an agent never does something is compatible with its having the power to do it.

\begin{proposition}\label{p:indep}
Let $\mathcal{M}$ be a model of the main argument, and $g$ a being maximally great at $@$ in $\mathcal{M}$. Then $\mathcal{M}$ has a causal extension in which $g$ causes every contingent being at every world accessible from $@$ at which that being exists, and a causal extension in which $g$ causes nothing. In both, $g$ is maximally great at $@$.
\end{proposition}

\begin{proof}
Set $C_1(x,y,w)$ iff $x = g$, $w$ is accessible from $@$, $y$ is contingent and $E(y,w)$; set $C_2$ empty. Both satisfy factivity, since $g$ exists at every world accessible from $@$. Neither extension alters $W$, $R$, $E$ or $\mathrm{ME}$, and by Assumption~\ref{a:powers} the extension of $\mathrm{ME}$ is not affected by which causal relation is added. Hence Assumptions~2.1 and~3.1 still hold, and $g$ is maximally great at $@$ by Theorem~3.4 of REL-01-a.
\end{proof}

\begin{corollary}\label{c:none}
No proposition about what the maximally great being causes, other than one that holds in every causal extension, follows from the premises of the main argument together with Assumption~\ref{a:powers}.
\end{corollary}

Between the two extremes of Proposition~\ref{p:indep} lie all intermediate patterns, and the same construction yields each. The main result is therefore compatible with a God who creates the world, with one on whom the world depends without being caused (the dependence of REL-01-d), and with one who does nothing.

\section{Consequences}

Four doctrines are thereby placed outside the reach of the main result. \emph{Creation}: that contingent beings exist because the maximally great being brings them about does not follow, since $C_2$ is a model. \emph{Providence}: that the course of the world is directed toward ends by that being does not follow, for the same reason. \emph{Miracle}: whether testimony could justify belief in a miracle is the question of REL-07, which presupposes a prior that does not exclude miracles; the main result neither raises nor lowers that prior. \emph{Prayer}: REL-12 argues that petitionary prayer makes a difference without causal efficacy; the main result neither supports nor threatens that account, since it says nothing about why outcomes occur.

\begin{objection}
Moral perfection requires action. A morally perfect being that could prevent suffering and does nothing is not morally perfect.
\end{objection}

\begin{reply}
On the definition of REL-01-a and REL-02-a, moral perfection excludes wrong options, and whether inaction is wrong depends on which goods are at stake. REL-02-a argues that some goods, those not brought about by God, are realizable only if God refrains from settling outcomes. Moral perfection therefore does not fix a pattern of action independently of facts about value, and the main result contains no such facts. A premise that the maximally great being must act in a given way would be an addition to the argument, of the kind that \citet{Swinburne2004} supplies by appeal to the probability that a perfectly good being would create.
\end{reply}

\begin{objection}
Omnipotence without exercise is an idle attribute, and an idle God is no God.
\end{objection}

\begin{reply}
The objection concerns the religious adequacy of the conclusion, not its truth. REL-01-a already disclaims identification with the God of any tradition. What the present note adds is that adequacy in the respect the objection names, a God who acts, requires premises about agency, and that those premises are not supplied by the modal argument or by its cosmological counterpart.
\end{reply}

\section{Conclusion}

The main result establishes the existence of a being with maximal powers, knowledge and goodness, and nothing about what that being does. Every claim about divine action must be argued for on other grounds, and the strength of the main result does not transfer to it.

\begin{thebibliography}{99}\small
\bibitem[Pike(1969)]{Pike1969} Pike, N. (1969). Omnipotence and God's ability to sin. \emph{American Philosophical Quarterly}, 6(3), 208--216.
\bibitem[Swinburne(2004)]{Swinburne2004} Swinburne, R. (2004). \emph{The Existence of God} (2nd ed.). Clarendon Press.
\end{thebibliography}

\end{document}
